\subsection{Independent terminal tests}

We now estimate the probability that all the local tests fail at one
fixed scale.  The selected leaves may depend on the selectors.  To keep
that dependence explicit, we first expose all selectors, identify the
terminal entries being read, and then average over the selectors.

\begin{lemma}[The fixed-scale estimate]
\label{lem:fixed-scale}
Fix a stable center \(x\) and an atom \(\xi\) of the center-exposure
sigma-field \(\mathcal F_x\) in \eqref{eq:center-exposure}, on
which the first default node \(U\) exists.  Let \(h\) be its height,
put \(r=r_h\), and use the local predicates \(Q_i\) for
\(e_i=(U,i)\), \(1\le i<M\), defined in \eqref{eq:local-predicate}, with \(y_n(t)=x+t2^{-n}\). For every fixed
\(t\in[1,2]\),
\begin{equation}
 \Pbb\left(
 Q_i(y_n(t))=0
 \text{ for all }1\le i<M,\ n\in W_{e_i}
 \,\middle|\,\xi
 \right)
 =(1-p/2)^{(M-1)r}.
 \label{eq:fixed-scale-miss}
\end{equation}
\end{lemma}

\begin{proof}
On the fixed atom \(\xi\), the node \(U\), the windows \(W_{e_i}\),
and the set of pairs
\[
 \mathcal I=\{(i,n):1\le i<M,\ n\in W_{e_i}\}
\]
are fixed.  Its cardinality is \((M-1)r\).  For these pairs define
\[
 A_{i,n}=S_{e_i}\bigl(J_{b_{e_i}}(y_n(t))\bigr).
\]
We first check that these are distinct selector entries, none of which
was exposed at \(x\).

Fix \(i\), and abbreviate \(a=a_{e_i}\) and \(b=b_{e_i}\).
For \(n<n'\) in this window,
\[
 y_n(t)-y_{n'}(t)
 =t\,2^{-n'}(2^{n'-n}-1)\ge 2^{-b}.
\]
Also \(y_n(t)-x=t2^{-n}\ge2^{-b}\).
All these points lie in \([x,x+2^{1-a}]\), an interval of length
at most \(1/4\), since \(a\ge3\).
Thus the circular distance between any two of them equals their
ordinary distance and is at least \(2^{-b}\).
Two points in the same half-open \(J_b\)-cell have circular distance
less than \(2^{-(b+2)}\).  Consequently the keys of
\(x\) and all \(y_n(t)\), \(n\in W_{e_i}\), are pairwise distinct.
This argument also covers a crossing of an integer and points on cell
boundaries.

For different \(i\), the selector tables themselves differ.
The variables \(A_{i,n}\) therefore read distinct coordinates outside
the entire exposed coordinate set defining \(\mathcal F_x\).
Conditional on \(\xi\), they remain independent fair bits.

Next condition on the sigma-algebra \(\mathcal S\) of all selectors.
For each pair, its selected local leaf
\[
 L_{i,n}=L_i(y_n(t))
\]
and its terminal address
\[
 \kappa_{i,n}=
 \bigl(L_{i,n},J_{b(L_{i,n})}(y_n(t))\bigr)
\]
are now fixed.  We claim that these addresses are distinct for all
pairs in \(\mathcal I\).
If the child indices differ, the leaves lie in disjoint child subtrees
of \(U\), so their terminal tables differ.
If the child indices agree but the leaves differ, their tables again
differ.  It remains to consider two different indices \(n,n'\) that
select the same leaf \(L\) below one child \(U_i\).

For every leaf below \(U_i\), its terminal resolution satisfies
\(b(L)\ge b_{e_i}\).  If \(U_i\) is a leaf, then \(b(L)=b_{e_i}\).
Otherwise the last edge into \(L\) occurs after \(e_i\) in the preorder,
so its window endpoint is larger.
By the nested-key identity~\eqref{eq:grid-nesting}, different keys at
resolution \(b_{e_i}\) remain different at resolution \(b(L)\),
including at grid boundaries.
The preceding separation of \(y_n(t)\) and \(y_{n'}(t)\) proves
\(\kappa_{i,n}\ne\kappa_{i,n'}\).
This proves the claim, even for pairs with \(A_{i,n}=0\).

The terminal tables are independent of \(\mathcal S\).
For each realization of all selectors, the addresses just identified
are fixed distinct coordinates of those tables.  The corresponding
terminal bits
\[
 V_{i,n}=
 T_{L_{i,n}}\bigl(J_{b(L_{i,n})}(y_n(t))\bigr)
\]
therefore have the product Bernoulli \(p\) law conditional on
\(\mathcal S\).  Since \(Q_i(y_n(t))=A_{i,n}V_{i,n}\),
only the pairs with \(A_{i,n}=1\) require a terminal bit to vanish.
On \(\xi\) we obtain
\[
 \Pbb\left(
 Q_i(y_n(t))=0\text{ for all }(i,n)\in\mathcal I
 \,\middle|\,\mathcal S
 \right)
 =(1-p)^{\sum_{(i,n)\in\mathcal I}A_{i,n}}.
\]
The empty-success case has exponent zero and probability one.

Finally, \(\mathcal F_x\subset\mathcal S\), so conditional expectation
and the independent fair law of the \(A_{i,n}\) on \(\xi\) give
\begin{align*}
 &\Pbb\left(
 Q_i(y_n(t))=0\text{ for all }(i,n)\in\mathcal I
 \,\middle|\,\xi
 \right)\\
 &\qquad=
 \E\left[(1-p)^{\sum_{(i,n)\in\mathcal I}A_{i,n}}
          \,\middle|\,\xi\right]\\
 &\qquad=
 \prod_{(i,n)\in\mathcal I}
 \left(\frac12+\frac12(1-p)\right)
 =(1-p/2)^{(M-1)r}.
\end{align*}
\end{proof}

The estimate applies separately to each scale fixed on the exposure
atom.  In the next step, the grid geometry will supply finitely many
such scales, allowing a union bound over all \(t\in[1,2]\).
